\section{Small information with an interior mean}\label{sec:small}
We now fix only $a_1<1$ and the correlation interval $[r_0,r_1]$.
Write the orthogonal decomposition
\[
 u=m+\ell+W,\qquad \ell=\sum_i c_iX_i,\quad
 s^2=\|\ell\|_2^2,\quad E_{\ge2}=\|W\|_2^2,
 \quad b=\max_i|c_i|.
\]
For a maximizing coordinate let
$B_0=\|u-bX_i\|_2^2=m^2+s^2-b^2+E_{\ge2}$.
Conditional Jensen and the convexity in $m$ of
$[\psi(m+b)+\psi(m-b)]/2-\psi(m)$ show $b\le a$.
The linear Bregman term at $bX_i$ cancels, so
\begin{equation}\label{eq:modal}
 B_0\le d_a=B_0+(a-b)^2\le(1+2U_{a_1})B_0,
 \qquad s^2+E_{\ge2}=\Var(u)\le2\psi(a)\le C a^2.
\end{equation}
Indeed $(a-b)^2\le2(\psi(a)-\psi(b))$ and
$\psi(a)-\psi(b)\le\E B_\psi(u,bX_i)\le U_{a_1}B_0$.

\begin{proposition}\label{prop:interiormean}
For every $m_0<1$ there exist $a_*,c_*>0$ such that
$|m|\le m_0$, $0<a\le a_*$, and $\cI(u)=\psi(a)$ imply
\[
 \Delta_r(u)\ge c_* a^4 B_0.
\]
The constants do not depend on dimension or an interior margin for $u$.
\end{proposition}
We prove it in three steps. All errors and constants in those steps
are uniform in dimension, $|m|\le m_0$, and the fixed correlation range.

\subsection{Control of the original higher modes}
For a first-chaos polynomial of norm $s$, the inequality
$\E e^{t\ell}\le e^{t^2s^2/2}$ follows from $\cosh x\le e^{x^2/2}$.
It gives $\|\ell\|_j\le C_js$ and, for fixed $\eta_0>0$,
\begin{equation}\label{eq:cliptail}
 \|\ell-C_{\eta_0}(\ell)\|_j\le C_j e^{-c_j/s^2},
\end{equation}
where $C_\eta$ clips to $[-\eta,\eta]$. At $s=0$ this error is zero.
Choose $\eta_0$ small relative to $1-m_0$ and put
$v=m+C_{\eta_0}(\ell)$, $t=\ell-C_{\eta_0}(\ell)$,
$e=u-v=W+t$. Symmetry gives $\E v=m$, $\E e=0$.
Both $v,T_rv$ lie in a fixed interior interval.
Let $E_a$ denote an error bounded by $Ce^{-c/a^2}$, with constants
allowed to change between occurrences.

The curve $f_r$ is concave by the scalar calculation above. Its tangent
at $\psi(a)$, with $k=rV(ra)/V(a)$, gives for any bounded comparison
field $v$ of amplitude $\beta=\psi^{-1}(\cI(v))$
\begin{equation}\label{eq:tangent}
 \Delta_r(u)\ge\Delta_r(v;\beta)
 +k[\cI(u)-\cI(v)]-[\cI(T_ru)-\cI(T_rv)].
\end{equation}
There is no ordering condition on the two input informations.
For the clipped comparison above the linear Bregman contribution is
\[
 C_{\rm lin}=k\E[V(v)e]-\E[V(T_rv)T_re],\qquad
 |C_{\rm lin}|\le Ca^2\sqrt{E_{\ge2}}+E_a.
\]
To see the bound, expand $V(m+z)=V(m)+z/(1-m^2)+O(z^2)$.
The constant terms vanish, and the first-chaos terms against $W$
vanish by orthogonality and self-adjointness. Moment bounds control the
quadratic remainder by $Ca^2\|W\|$, and \eqref{eq:cliptail} handles $t$.

We claim that for small $a$
\begin{equation}\label{eq:entrancepayment}
 k\E B_\psi(u,v)-\E B_\psi(T_ru,T_rv)
 \ge cE_{\ge2}-E_a.
\end{equation}
Here are details retaining large input residuals. Split $e=G+Z$ at a
small fixed threshold $\eta_1$. The bad support mass satisfies
$P\le Ca^2/\eta_1^2+E_a$. On the good set the input quadratic
coefficient can be made arbitrarily close to $1/[2(1-m^2)]$ by choosing
$\eta_0,\eta_1$ small. On the bad set keep the global payment
$k\|Z\|^2/2$. The output has the same approximate quadratic coefficient
and a remainder at most $C\|e\|^p$ by \eqref{eq:hyper}.
Since $k=r^2+O(a^2)$, the constant-coefficient quadratic form has a
uniform positive gap on degrees at least two:
$(r^2-r^4)/[2(1-m^2)]$.
Its low-degree block has bounded norm. By \eqref{eq:lowproj},
$\|\Pi_{\le1}G\|^2\le CP^{1/2}\|Z\|^2+E_a$, because
$\Pi_{\le1}e=\Pi_{\le1}t$.
The bad output square is bounded by $CP^{2\theta}\|Z\|^2$,
and its cross term by $CP^\theta\|G\|\|Z\|$.
Young's inequality absorbs this into a fixed fraction of the good gap
and $CP^{2\theta}\|Z\|^2$.
Finally $\|e\|^p\le Ca^{p-2}(\|G\|^2+\|Z\|^2)+E_a$.
As $a\to0$, every bad coefficient tends to zero and can be absorbed
into the retained input payment. This proves \eqref{eq:entrancepayment}.
Use $\Delta_r(v;\beta)\ge0$ in \eqref{eq:tangent}. We obtain
\begin{equation}\label{eq:entrance}
 \Delta_r(u)\ge cE_{\ge2}-Ca^2\sqrt{E_{\ge2}}-E_a,
 \qquad E_{\ge2}\le C(\Delta_r(u)+a^4).
\end{equation}

\subsection{A quartic expansion retaining bad entropy}
Assume now $E_{\ge2}\le K_0a^4$, with fixed $K_0$.
Retain the preceding $v,e,t$ and split $e=G+Z$ at the shrinking threshold
$\tau=a^{1/2}$. Its bad mass is $P\le C_{K_0}a^3+E_a$.
Set $h=1-m^2$ and
\[
 A_{\rm bad}=\E[B_\psi(u,v)\mathbf1_{\{|e|>\tau\}}].
\]
Then
\[
 \tfrac12\|Z\|^2\le A_{\rm bad}\le U_\gamma\|Z\|^2,
 \quad \|T_rZ\|=o(a^2),\quad \|\Pi_{\le1}G\|=o(a^2),
 \quad |\E\ell^2Z|\le Ca^2P^{1/4}\|Z\|=o(a^4).
\]
These follow from \eqref{eq:bregman}, \eqref{eq:hyper},
\eqref{eq:lowproj}, and fourth/eighth moment bounds for $\ell$.
Taylor expansion only on the good residual set gives
$\E[B_\psi(u,v)\mathbf1_{\rm good}]=\|G\|^2/(2h)+o(a^4)$.
For its multiplier, first restrict to $|\ell|\le a^{1/2}$; the variation
is $O(a^{1/2})$. The complement has mass exponentially small in $1/a$,
and $|G|\le\tau$ controls its contribution. The cubic residual cost
is $O(\tau\|G\|^2)$.
At the output the remainder is $O(a^{2p})=o(a^4)$, while H\"older
bounds multiplier variation by
$C\|T_rv-m\|_{p/(p-2)}\|T_re\|_p^2=O(a^5)+E_a$.
The first-chaos linear terms cancel for the full $e=W+t$ before the
bad part is removed.

Writing $C_4=(1+3m^2)/(12h^3)$, these estimates give
\begin{align}
 \psi(a)&=\frac{s^2}{2h}+C_4\E\ell^4+\frac m{h^2}\E\ell^2G
       +\frac{\|G\|^2}{2h}+A_{\rm bad}+o(a^4),\label{eq:quarticin}\\
 \cI(T_ru)&=\frac{r^2s^2}{2h}+C_4r^4\E\ell^4
       +\frac{mr^2}{h^2}\E[\ell^2T_rG]
       +\frac{\|T_rG\|^2}{2h}+o(a^4).\label{eq:quarticout}
\end{align}
For example, expanding $V(m+\ell)$ through order two leaves an
$L^2$ remainder $O(a^3)$ paired with an $O(a^2)$ residual.
Equation~\eqref{eq:cliptail} pays all clipping terms, and the preceding
bound on $\E\ell^2Z$ removes the bad cubic coupling. This proves
uniformity of these errors without a pointwise assumption on $W$.

Equation~\eqref{eq:quarticin} gives $s^2=ha^2+O(a^4)$.
For $\kappa=\sum_i c_i^4/s^4$ and $G_j=\Pi_jG$, eliminate $s^2/(2h)$
and complete the level-two square to obtain
\begin{align}
 \Delta_r(u)={}&r^2(1-r^2)a^4
       \left\{\frac{\kappa m^2}{3h}+\frac{1-\kappa}{6}\right\}
   +\frac{r^2(1-r^2)}{2h}
       \left\|G_2+\frac m h(\ell^2-s^2)\right\|^2\notag\\
 &+\sum_{j\ge3}\frac{r^2-r^{2j}}{2h}\|G_j\|^2
       +r^2A_{\rm bad}+o(a^4).\label{eq:normal}
\end{align}
Indeed $\E\ell^4=s^4+2(s^4-\sum c_i^4)$, and square completion
subtracts $r^2(1-r^2)m^2\,2(s^4-\sum c_i^4)/(2h^3)$.
Substitute $s^4=h^2a^4+O(a^6)$ and
$\psi(ra)-r^2\psi(a)=-r^2(1-r^2)a^4/12+O(a^6)$ to verify the
bracket. It is at least $(m^2+1-\kappa)/12$.
Hence $\Delta=o(a^4)$ forces
\begin{equation}\label{eq:coneentrance}
 m\to0,\quad\kappa\to1,\quad E_{\ge2}=o(a^4),\quad
 s^2-b^2=o(a^2),\quad b^2/(ha^2)\to1.
\end{equation}
For the higher energy use $A_{\rm bad}\ge\|Z\|^2/2$, the negligible
low projection of $G$, and
$\|(m/h)(\ell^2-s^2)\|=o(a^2)$ in \eqref{eq:normal}.
Also $\kappa\le b^2/s^2$. The absolute error $o(a^4)$ is used only
for this entrance statement.

\subsection{Relative coercivity near a coordinate}
There are small fixed $\mu,\varepsilon,\delta,a_c>0$ such that, if
\begin{equation}\label{eq:cone}
 |m|\le\mu,\quad \ell_0=\sum_{i\ne1}c_iX_i,\quad
 l_0=\|\ell_0\|\le\varepsilon a,\quad
 |b^2/((1-m^2)a^2)-1|\le\delta,\quad a\le a_c,
\end{equation}
with $c_1=b\ge0$, then
\begin{equation}\label{eq:relative}
 \Delta_r(u)\ge c(a^4m^2+a^2l_0^2+E_{\ge2}).
\end{equation}
To prove it, take the exact conditional field $v=m+bX_1$ and
$A=v+C_\eta(\ell_0)$ for fixed small $\eta$.
Put $t=\ell_0-C_\eta(\ell_0)$, $e=u-A=W+t$.
Choose the constants so that $|A|,|T_rA|\le\gamma$ with
$\gamma^2<1-r_1^2$.
All fixed-norm clipping errors are bounded by
$E_l=Ce^{-c/l_0^2}$, interpreted as zero when $l_0=0$.
Symmetry and independence of $X_1$ cancel the linear $v$-to-$A$ terms.
Equation~\eqref{eq:tangent} becomes
\begin{align}
 \Delta_r(u)\ge{}&\Delta_r(v;\beta)
 +k\E B_\psi(A,v)-\E B_\psi(T_rA,T_rv)\notag\\
 &+k\E[V(A)e]-\E[V(T_rA)T_re]\notag\\
 &+k\E B_\psi(u,A)-\E B_\psi(T_ru,T_rA),\label{eq:conedecomp}
\end{align}
where $\beta=\psi^{-1}(\cI(v))$ and $k=f_r'(\psi(a))$.

For this genuine one-bit field, capacity and convexity of
$F(J)-(2/3)J^2$ give
$D(v)-F(\cI(v))\ge(2/3)m^2\cI(v)^2$.
Along its noise flow, \eqref{eq:budget} applies.
If the terminal gap is at least half $\psi(r\beta)$ it bounds
$c m^2\beta^4$ immediately. Otherwise its information at every earlier
time is at least its terminal value, at least $c\beta^2$;
integrating the production surplus gives the same bound.
From \eqref{eq:cone}, $\beta\asymp a$. Thus
\begin{equation}\label{eq:onebitpay}
 \Delta_r(v;\beta)\ge c m^2a^4.
\end{equation}

The cubic terms in the next line of \eqref{eq:conedecomp}'s first row
average to zero, since the clipped first chaos is symmetric and independent
of $X_1$. Its fourth moments are $O(l_0^4)$. Consequently this contribution
is $(Q_{11}/2)l_0^2+O(l_0^4)+O(E_l)$, where
\[
 Q_{11}=k\E(1-v^2)^{-1}-r^2\E(1-(T_rv)^2)^{-1}
 =r^2(1-r^2)\left\{\frac{(1+3m^2)b^2}{(1-m^2)^3}
                 -\frac{a^2}{3(1-m^2)}\right\}+O(a^4).
\]
Here $k=r^2-r^2(1-r^2)a^2/3+O(a^4)$; all expansions are analytic
on the fixed interior interval. At $m=0$, $b^2=a^2$, the leading
coefficient is $2r^2(1-r^2)/3$. Choose $\mu,\delta,a_c$ small and
then $\varepsilon$ small to obtain
\begin{equation}\label{eq:otherpay}
 k\E B_\psi(A,v)-\E B_\psi(T_rA,T_rv)
 \ge ca^2l_0^2-CE_l.
\end{equation}

Strong convexity at the input, \eqref{eq:bregman} at the output,
and the degree-at-least-two spectral gap give for the last row
\begin{align}
 k\E B_\psi(u,A)-\E B_\psi(T_ru,T_rA)
 &\ge\left\{\frac k2-\frac{r^4}{2(1-\gamma^2)}\right\}E_{\ge2}
       -C E_{\ge2}^{p/2}-CE_l\notag\\
 &\ge cE_{\ge2}-CE_l.\label{eq:higherpay}
\end{align}
In the first line a fixed arbitrarily small fraction of the positive
coefficient may be spent on cross terms involving $t$; its exponentially
small norm is included in $CE_l$. The coefficient remains positive by
our choice of $\gamma$, and $E_{\ge2}\le Ca^2$ absorbs the remainder.

Finally expand $V(v+\ell_0)$ through order two in $\ell_0$ while
keeping its full dependence on $v$. Conditional orthogonality removes
$V(v)W$ and the even-in-$X_1$ part of $\psi''(v)\ell_0 W$.
The odd coefficient is at most $C|m|b$, and
$|\psi'''(v)|\le C(|m|+b)$. The cubic remainder has norm $O(l_0^3)$.
The identical calculation after noise, using self-adjointness, gives
\begin{equation}\label{eq:coupling}
 |k\E[V(A)e]-\E[V(T_rA)T_re]|
 \le C\{|m|a l_0+(|m|+a)l_0^2+l_0^3\}\sqrt{E_{\ge2}}+CE_l.
\end{equation}
Young's inequality spends half of \eqref{eq:higherpay}; its other cost
is bounded by
$C(\mu^2+\varepsilon^2+\mu^2\varepsilon^2+\varepsilon^4)a^2l_0^2$,
which is absorbed in \eqref{eq:otherpay} by decreasing $\mu,\varepsilon$.
For $l_0\le\varepsilon a$,
$E_l/(a^2l_0^2)\le C\varepsilon^2e^{-c/l_0^2}/l_0^4\to0$
uniformly as $a\to0$; the error is exactly zero when $l_0=0$.
Thus the clipping costs are relative even if $B_0$ is exceptionally
small. Combining \eqref{eq:conedecomp}--\eqref{eq:coupling} proves
\eqref{eq:relative}.

\begin{proof}[Proof of Proposition~\ref{prop:interiormean}]
Failure would give $a_j\to0$ and $B_{0,j}>0$ with
$\Delta_j/(a_j^4B_{0,j})\to0$. Since $B_{0,j}\le m_0^2+Ca_j^2$,
we have $\Delta_j=o(a_j^4)$. Equation~\eqref{eq:entrance} gives
$E_{\ge2}=O(a_j^4)$, and \eqref{eq:normal} implies
\eqref{eq:coneentrance}. Choose the largest singleton and orient it.
The fixed conditions \eqref{eq:cone} eventually hold. Then
\eqref{eq:relative} gives
$\Delta_j\ge c(a_j^4m_j^2+a_j^2l_{0,j}^2+E_{\ge2,j})
\ge ca_j^4B_{0,j}$, a contradiction. If $B_0=0$, \eqref{eq:modal}
already gives $d_a=0$.
\end{proof}
